% 大学物理核心公式与考点汇总 - XeLaTeX无错优化版（支持五级标题TOC）
% 层级：chapter(1) > section(2) > subsection(3) > subsubsection(4) > subsubsubsection(5)
% 编号示例：1.1.1.1.1，适配MiKTeX/TeXLive Windows
\documentclass[12pt,a4paper,oneside]{report}
% ======================== 核心宏包 ========================
\usepackage{geometry}
\usepackage{microtype}  % 优化字符间距，减少Underfull/Overfull警告
\usepackage{needspace}
% 中文与字体【修复顺序冲突】
\usepackage[UTF8, fontset=windows]{ctex}
\usepackage{fontspec}
\setmainfont{Times New Roman}
\setsansfont{Microsoft YaHei}
% 数学与格式
\usepackage{amsmath,amssymb,amsfonts}
\usepackage{upgreek}
\usepackage{xcolor}
\usepackage{colortbl}
\usepackage{booktabs}
\usepackage{multirow}
\usepackage{makecell}
\usepackage{siunitx}
\usepackage{titlesec}
\usepackage{titletoc}
\usepackage{tcolorbox}
\usepackage{fancyhdr}
\usepackage{enumitem}
\usepackage{tikz}
\usetikzlibrary{decorations.pathmorphing} % 修复snake波浪线报错
% 超链接顺序规范【hypcap必须在hyperref之后加载】
\usepackage[
    bookmarks=true,bookmarksdepth=5,unicode=true,colorlinks=true,
    linkcolor=mainblue,citecolor=mainblue,urlcolor=mainblue,
    pdfborder={0 0 0}
]{hyperref}
\usepackage{hypcap}
% 告诉 hyperref 五级标题的目录层级，消除 bookmark 警告
\makeatletter
\providecommand\toclevel@subsubsubsection{5}
\makeatother
% ======================== 基础页面设置 ========================
\geometry{a4paper,left=2.5cm,right=2.5cm,top=2.5cm,bottom=2.5cm,headheight=15pt}
% 全局颜色定义
\definecolor{mainblue}{RGB}{41, 90, 141}
\definecolor{secondcyan}{RGB}{26, 117, 117}
\definecolor{thirdteal}{RGB}{27, 120, 101}
\definecolor{highlight}{RGB}{255, 240, 200}
\definecolor{tablehead}{RGB}{230, 240, 250}
\definecolor{textgray}{RGB}{60, 60, 60}
% 页眉页脚全局样式
\pagestyle{fancy}
\fancyhf{}
\fancyhead[L]{\color{mainblue}\bfseries\small 大学物理核心公式与考点汇总}
\fancyhead[R]{\color{mainblue}\thepage}
\renewcommand{\headrulewidth}{0.5pt}
\renewcommand{\headrule}{\color{mainblue}\hrule width\headwidth height\headrulewidth}
% 第一章首页plain页眉统一
\fancypagestyle{plain}{
    \fancyhf{}
    \fancyhead[L]{\color{mainblue}\bfseries\small 大学物理核心公式与考点汇总}
    \fancyhead[R]{\color{mainblue}\thepage}
    \renewcommand{\headrulewidth}{0.5pt}
    \renewcommand{\footrulewidth}{0pt}
}
% 各级标题垂直间距紧凑优化（新增五级标题间距，titlesec接管全部间距）
\titlespacing*{\section}{0pt}{6pt plus 2pt minus 2pt}{3pt plus 2pt minus 2pt}
\titlespacing*{\subsection}{0pt}{4pt plus 2pt minus 2pt}{2pt plus 2pt minus 2pt}
\titlespacing*{\subsubsection}{0pt}{3pt plus 1pt minus 1pt}{1pt plus 1pt minus 1pt}
\titlespacing*{\subsubsubsection}{0pt}{2pt plus 1pt minus 1pt}{1pt plus 1pt minus 1pt}
% 枚举列表全局紧凑配置
\setlist{
    itemsep=0.1ex,
    parsep=0.2ex,
    leftmargin=1.5em,
    labelsep=0.4em,
    topsep=0.2ex,
    partopsep=0.1ex,
    font=\small
}
\setlist[itemize,2]{leftmargin=1.8em,label=$\cdot$,itemsep=0.05ex,parsep=0.1ex}
% 数学公式上下间距压缩
\numberwithin{equation}{section}
\allowdisplaybreaks
\setlength{\abovedisplayskip}{3pt}
\setlength{\belowdisplayskip}{3pt}
\setlength{\abovedisplayshortskip}{2pt}
\setlength{\belowdisplayshortskip}{2pt}
\makeatletter
\renewcommand{\@mathmargin}{0pt}
\makeatother
% 整体排版容错、缩进、段落间距
\setlength{\emergencystretch}{2em}
\sloppy
\setlength{\parindent}{0pt}
\setlength{\parskip}{0.8ex}
\setlength{\tabcolsep}{4pt}
% ======================== 【缺失定义，子文件依赖，新增】 ========================
% solution 解题环境，子文件大量 \begin{solution}
\newenvironment{solution}{\par\noindent\textbf{解：}\ignorespaces}{\par}
% \vect 矢量命令，子文件大量 \vect{r} \vect{i}
\newcommand{\vect}[1]{\vec{#1}}
% ======================== 自定义快捷命令 ========================
% 高亮核心公式框
\newcommand{\keyformula}[1]{%
    \needspace{2\baselineskip}%
    \begin{tcolorbox}[
        colback=highlight,colframe=white,boxrule=0pt,arc=0pt,
        left=6pt,right=6pt,top=4pt,bottom=4pt,center upper,
        boxsep=0pt
    ]
    $\displaystyle #1$
    \end{tcolorbox}%
    \vspace{0.2em}
}
% 符号说明命令
\newcommand{\symb}[2]{\textbf{#1}：\parbox[t]{10cm}{\raggedright #2}\\[0.3ex]}
% 核心考点蓝色边框提示框
\newcommand{\corepoint}[1]{%
    \begin{tcolorbox}[
        colback=white,colframe=mainblue,boxrule=0.5pt,arc=2pt,
        left=6pt,right=6pt,top=2pt,bottom=2pt,center upper,
        boxsep=0pt
    ]
    \textbf{\color{mainblue}核心考点：#1}
    \end{tcolorbox}%
    \vspace{0.2em}
}
% ======================== 【核心：完整定义五级标题 subsubsubsection + TOC目录渲染】 ========================
\makeatletter
% 1. 创建五级计数器，依附于四级 subsubsection，新建计数器默认=0，\refstepcounter后变为1，不会输出0
\newcounter{subsubsubsection}[subsubsection]
% 2. 五级编号格式：1.1.1.1.1 链式编号
\renewcommand{\thesubsubsubsection}{\thesubsubsection.\arabic{subsubsubsection}}
% 3. 定义五级标题命令 \subsubsubsection{标题名称}，层级编号=5
\newcommand{\subsubsubsection}{\@startsection
    {subsubsubsection}
    {5}
    {\z@}
    {3.25ex \@plus1ex \@minus.2ex}
    {1.5ex \@plus.2ex}
    {\normalfont\normalsize\bfseries}
}
% 4. 重定义目录内部命令 \l@subsubsubsection，让tocdepth=5时正常生成五级目录条目
\def\l@subsubsubsection#1#2{%
  \ifnum \c@tocdepth >4  % tocdepth>4 才渲染五级
    \addpenalty{-\@highpenalty}%
    \vskip 0.4em \@plus\p@
    \setlength\@tempdima{16em}% 五级目录缩进宽度
    \begingroup
      \parindent \z@ \rightskip \@pnumwidth
      \parfillskip -\@pnumwidth
      \leavevmode
      \advance\leftskip\@tempdima
      \hskip -\leftskip
      #1\nobreak\hfil \nobreak\hb@xt@\@pnumwidth{\hss #2}\par
      \penalty\@highpenalty
    \endgroup
  \fi
}
\makeatother
% ======================== 目录深度关键配置（开启五级显示） ========================
\setcounter{secnumdepth}{5}   % 正文允许编号到第5级标题
\setcounter{tocdepth}{5}     % 目录TOC渲染显示到第5级标题
% ======================== 文档入口 ========================
\begin{document}
% 封面页
% 2.cover.tex 封面独立子文件
% 被 main.tex 调用

\newcommand{\makeCover}
{
  \begin{titlepage}
    \newgeometry{left=2cm, right=2cm, top=2cm, bottom=2cm}

    \centering
    \vspace*{1cm}

    {\Large \itshape 以习题夯实物理逻辑，\\以推导巩固数理思维，\\梳理力学、电磁学、振动波动核心题型}\\[0.5cm]
    \rule{\textwidth}{1pt}\\[0.5cm]

    {\Huge \bfseries \color{mainblue} 大学物理}\\[8pt]
    {\Huge \bfseries \color{mainblue} 课后作业习题集}\\[12pt]
    {\Large \itshape —— 分模块刷题与完整解题步骤 ——}\\[1cm]

    \rule{0.6\textwidth}{0.5pt}\\[1cm]

    {\Large \textbf{习题模块：} 力学 | 电磁学 | 振动\&波动}\\[8pt]
    {\Large \textbf{文档用途：} 课后作业完成 · 题型归纳 · 考前刷题}\[1.5cm]

    {\Large \textbf{整理者：}浮萍}\\[6pt]
    {\large \textbf{联系邮箱：}\href{mailto:3071974740@qq.com}{3071974740@qq.com}}\\[6pt]
    {\large \textbf{GitHub：}\url{https://github.com/Duckweed-yhb}}\\[6pt]
    {\Large 整理日期：\today}\\[1cm]

    \rule{0.6\textwidth}{0.5pt}\\[1cm]

    {\large \textbf{版本：} v1.0 作业汇编版}\\[6pt]
    {\small \itshape Copyright © \the\year\ 浮萍 · 禁止商用 · 学习交流使用}\\[0.5cm]

    \rule{\textwidth}{1pt}
    \vfill

    \restoregeometry
  \end{titlepage}
}
\clearpage
\thispagestyle{fancy}
% 自动生成目录（两次XeLaTeX编译生效五级目录）
\tableofcontents
\clearpage
% ========== 第一篇 力学 ==========
\part{第一篇 力学}
\section{质点运动学习题解答}
\subsection{直角坐标分量求轨道、速度、加速度}
已知质点的运动方程
\[
x=\sqrt{3}\cos \frac{\pi}{4}t,\quad y=\sin \frac{\pi}{4}t
\]
式中，$x,y$ 以 $\mathrm{m}$ 计，$t$ 以 $\mathrm{s}$ 计。
\begin{enumerate}
    \item 求质点的轨道方程，并在 $xOy$ 平面上描画出质点的轨道；
    \item 求出质点的速度和加速度表示式；
    \item 求 $t=1\ \mathrm{s}$ 时质点的位置、速度和加速度。
\end{enumerate}
\begin{solution}
\begin{enumerate}
    \item \textbf{轨道方程}
由运动方程变形：
\[
\frac{x}{\sqrt{3}} = \cos\frac{\pi}{4}t,\quad y = \sin\frac{\pi}{4}t
\]
利用三角恒等式 $\cos^2\theta+\sin^2\theta=1$：
\[
\boxed{\frac{x^2}{3} + y^2 = 1}
\]
轨道为长半轴 $a=\sqrt{3}$（$x$ 轴）、短半轴 $b=1$（$y$ 轴）的椭圆。
    \item \textbf{速度与加速度表达式}
位置矢量：
\[
\vect{r} = \sqrt{3}\cos\frac{\pi}{4}t\,\vect{i} + \sin\frac{\pi}{4}t\,\vect{j}
\]
速度 $\vect{v}=\dfrac{d\vect{r}}{dt}$：
\[
\begin{aligned}
v_x &= \frac{dx}{dt} = -\frac{\sqrt{3}\pi}{4}\sin\frac{\pi}{4}t,\\
v_y &= \frac{dy}{dt} = \frac{\pi}{4}\cos\frac{\pi}{4}t,\\
\vect{v} &= -\frac{\sqrt{3}\pi}{4}\sin\frac{\pi}{4}t\,\vect{i} + \frac{\pi}{4}\cos\frac{\pi}{4}t\,\vect{j}.
\end{aligned}
\]
加速度 $\vect{a}=\dfrac{d\vect{v}}{dt}$：
\[
\begin{aligned}
a_x &= -\frac{\sqrt{3}\pi^2}{16}\cos\frac{\pi}{4}t = -\frac{\pi^2}{16}x,\\
a_y &= -\frac{\pi^2}{16}\sin\frac{\pi}{4}t = -\frac{\pi^2}{16}y,\\
\vect{a} &= -\frac{\pi^2}{16}\left(\sqrt{3}\cos\frac{\pi}{4}t\,\vect{i}+\sin\frac{\pi}{4}t\,\vect{j}\right)
= -\frac{\pi^2}{16}\vect{r}.
\end{aligned}
\]
    \item \textbf{$t=1\ \mathrm{s}$ 瞬时量}
$\dfrac{\pi}{4}t\big|_{t=1}=\dfrac{\pi}{4}$
\[
\begin{aligned}
&\vect{r}\big|_{t=1} = \sqrt{3}\cos\frac{\pi}{4}\,\vect{i}+\sin\frac{\pi}{4}\,\vect{j}
= \frac{\sqrt{6}}{2}\vect{i} + \frac{\sqrt{2}}{2}\vect{j}\ \mathrm{(m)},\\
&\vect{v}\big|_{t=1} = -\frac{\sqrt{6}\pi}{8}\vect{i}+\frac{\sqrt{2}\pi}{8}\vect{j}\ \mathrm{(m/s)},\\
&\vect{a}\big|_{t=1} = -\frac{\sqrt{6}\pi^2}{32}\vect{i}-\frac{\sqrt{2}\pi^2}{32}\vect{j}\ \mathrm{(m/s^2)}.
\end{aligned}
\]
\end{enumerate}
\end{solution}
\subsection{\texorpdfstring{变加速度 $a=-kv^2$ 积分求速度}{变加速度a=-kv^2积分求速度}}
一艘正在以 $v_0$ 速率沿直线行驶的电艇，在发动机关闭后，其加速度方向与速度方向相反，大小与速度平方成正比，即 $a=-kv^2$，式中 $k$ 为常量。试求电艇在关闭发动机后又行驶 $x$ 距离时的速度。
\begin{solution}
加速度作变量代换 $a=\dfrac{dv}{dt}=\dfrac{dv}{dx}\cdot\dfrac{dx}{dt}=v\dfrac{dv}{dx}$，代入：
\[
v\frac{dv}{dx} = -kv^2
\]
分离变量（$v\neq0$）：
\[
\frac{dv}{v} = -k\,dx
\]
定积分上下限：初始 $x=0,\,v=v_0$；末态 $x=x,\,v=v$
\[
\int_{v_0}^{v} \frac{dv}{v} = -k\int_{0}^{x}dx
\]
\[
\ln v - \ln v_0 = -kx \implies \ln\frac{v}{v_0}=-kx
\]
\[
\boxed{v = v_0 e^{-kx}}
\]
\end{solution}
\subsection{拉船靠岸的速度、加速度（牵连运动）}
\begin{center}
% 加载波浪线绘图库
\begin{tikzpicture}[scale=1.2, decoration={snake, amplitude=2pt, segment length=6pt}]
% 岸壁
\draw (0,0) -- (0,3);
\draw (-1,3) -- (0,3);
\fill (0,3) circle (2pt) node[above right]{收绳点};
% 高度h标注
\draw[<->] (-0.3,0) -- (-0.3,3) node[midway,left]{$h$};
% 水面波浪线（使用decorate）
\draw[decorate] (0,0) -- (4,0);
% 船位置
\fill (4,0) circle (2pt);
% 绳子
\draw (0,3) -- (4,0);
% 水平距离s
\draw[<->] (0,-0.3) -- (4,-0.3) node[midway,below]{$s$};
% 收绳速度v0
\draw[<-] (-0.8,3) -- (-0.2,3) node[above]{$v_0$};
\end{tikzpicture}
\end{center}
在离水面高度为 $h$ 的岸壁上，有人用绳子拉船靠岸，船位于离岸壁 $s$ 距离处，当以 $v_0$ 的速度收绳时，试问船的速度、加速度的大小各为多少？
\begin{solution}
设绳长为 $l$，几何关系：$l=\sqrt{h^2+s^2}$
收绳速率：$\dfrac{dl}{dt}=-v_0$（绳长缩短，导数为负）
船的位置坐标为 $s$，船速 $v=\dfrac{ds}{dt}$
两边对 $t$ 求导：
\[
\frac{dl}{dt} = \frac{s}{\sqrt{h^2+s^2}} \cdot \frac{ds}{dt}
\]
\[
-v_0 = \frac{s}{l}\,v \implies v = -\frac{l}{s}v_0 = -\frac{\sqrt{h^2+s^2}}{s}\,v_0
\]
船速大小（指向岸，速率）：
\[
\boxed{|v| = \frac{\sqrt{h^2+s^2}}{s}\,v_0}
\]
加速度 $a=\dfrac{dv}{dt}$：
\[
v = -v_0 (h^2+s^2)^{\frac12} s^{-1}
\]
\[
\frac{dv}{dt} = -v_0 \cdot \frac{s\cdot \dfrac{s}{\sqrt{h^2+s^2}} - \sqrt{h^2+s^2}}{s^2} \cdot \frac{ds}{dt}
\]
\[
\frac{dv}{dt} = -v_0 \cdot \frac{s^2-(h^2+s^2)}{s^2\sqrt{h^2+s^2}} \cdot v
= -v_0 \cdot \frac{-h^2}{s^2\sqrt{h^2+s^2}} \cdot \left(-\frac{\sqrt{h^2+s^2}}{s}v_0\right)
\]
\[
\boxed{a = -\frac{h^2 v_0^2}{s^3}}
\]
负号表示加速度指向河岸，加速度大小：$\boldsymbol{\dfrac{h^2 v_0^2}{s^3}}$。
\end{solution}
\subsection{鼓轮匀角加速转动，求C点速度加速度}
\begin{center}
\begin{tikzpicture}[scale=1]
\draw (-3,2) node[circle,fill,inner sep=2pt]{} node[above]{$A$} -- (0,2) node[circle,fill,inner sep=2pt]{} node[above]{$B$};
\draw (0,2) circle (1);
\draw (0,2) -- (0,1) node[circle,fill,inner sep=2pt]{} node[right]{$O$} node[midway,right]{$R$};
\draw (0,0) node[circle,fill,inner sep=2pt]{} node[below]{$C$};
\end{tikzpicture}
\end{center}
一卷扬机的鼓轮自静止开始作匀角加速度转动，水平绞索上的 $A$ 点经 $3\ \mathrm{s}$ 后到达鼓轮边缘上的 $B$ 点处。已知 $AB=0.45\ \mathrm{m}$，鼓轮半径 $R=0.5\ \mathrm{m}$。求 $A$ 点到达最低点 $C$ 时的速度与加速度。
\begin{solution}
\textbf{修正说明：A点随绞索先水平匀加速直线运动，缠绕上鼓轮后做匀角加速圆周运动，切向加速度保持不变。}
1. 水平段 $AB$：初速度为0匀加速直线运动
\[
s_{AB} = \frac12 a_t t_1^2 \implies 0.45 = \frac12 a_t \cdot 3^2
\]
\[
a_t = \frac{2\times 0.45}{9} = 0.1\ \mathrm{m/s^2}
\]
2. 鼓轮匀角加速关系：$a_t = R \beta \implies \beta = \dfrac{a_t}{R} = \dfrac{0.1}{0.5}=0.2\ \mathrm{rad/s^2}$
$A$ 由B转到最低点C，转过圆心角 $\theta=\pi$，匀角加速转动：
\[
\omega_C^2 - \omega_B^2 = 2\beta \pi
\]
先求B点角速度：$\omega_B = \dfrac{v_B}{R},\ v_B = a_t t_1 = 0.1\times3 = 0.3\ \mathrm{m/s}$
\[
\omega_B = \frac{0.3}{0.5}=0.6\ \mathrm{rad/s}
\]
\[
\omega_C^2 = 0.6^2 + 2\times 0.2 \times \pi = 0.36 + 0.4\pi
\]
3. C点线速度：
\[
v_C = R \omega_C = 0.5 \sqrt{0.36+0.4\pi} \approx 0.5 \times 1.351 = \boxed{0.676\ \mathrm{m/s}}
\]
4. 加速度分解：
切向加速度全程不变：$a_t=0.1\ \mathrm{m/s^2}$；
法向（向心）加速度：
\[
a_n = \frac{v_C^2}{R} = \frac{R^2\omega_C^2}{R}=R\omega_C^2 = 0.5(0.36+0.4\pi)\approx 0.456\ \mathrm{m/s^2}
\]
合加速度大小：
\[
a = \sqrt{a_t^2+a_n^2}
= \sqrt{0.1^2 + 0.456^2}
\approx \boxed{0.467\ \mathrm{m/s^2}}
\]
\end{solution}
\subsection{相对运动求风速}
有人以 $v=3\ \mathrm{m\cdot s^{-1}}$ 的速率向东奔跑，他感到风从北方吹来；当奔跑的速率加倍时，则感到风从东北方向吹来，求风的速度。
\begin{solution}
建立坐标系：东为 $x$ 正方向，北为 $y$ 正方向。设风速 $\vect{u}=u_x \vect{i}+u_y\vect{j}$。
人的速度 $\vect{v}_\text{人}$，相对风速 $\vect{u}_\text{相对}=\vect{u}-\vect{v}_\text{人}$。
\begin{enumerate}
    \item $v_\text{人}=3\vect{i}$，相对风正北吹来，相对速度无x分量：
    \[
    u_x - 3 = 0 \implies u_x=3\ \mathrm{m/s}
    \]
    \item 速率加倍 $v_\text{人}=6\vect{i}$，风从东北吹来（相对速度指向西南：$u_x-6 <0,\ u_y<0$，x,y分量相等）：
    \[
    u_x-6 = u_y \implies 3-6=u_y \implies u_y=-3\ \mathrm{m/s}
    \]
\end{enumerate}
风速矢量：
\[
\vect{u}=3\vect{i}-3\vect{j}\ \mathrm{(m/s)}
\]
大小：
\[
u=\sqrt{3^2+(-3)^2}=3\sqrt{2}\ \mathrm{m/s}
\]
方向：\textbf{风源自西北方向（西北风，吹向东南）}。
\end{solution}
\subsection{流速与离岸距离成正比的小船渡河轨迹【重点修复】}
河宽为 $l$，水流速与离岸的距离成正比，中心流速最大为 $v_0$，两岸处流速为零。一小船以恒定的相对速度 $v$ 垂直于水流从一岸驶向另一岸。当它驶至河宽的 $1/4$ 处时，发现燃料不足，立即掉头以相对速度 $0.5v$ 垂直于水流驶回原岸，求此船驶往对岸的轨迹及返回到原岸的地点。
\begin{solution}
设垂直河岸为 $y$ 轴（$0\le y\le l$），沿河水流向为 $x$ 轴。
流速分布：$u(y) = \dfrac{2v_0}{l}y$（线性分布，$y=l/2$ 时 $u=v_0$）。
\subsubsection{第一阶段：去程 \texorpdfstring{$0\to l/4$}{0到l/4}，船垂直对岸 \texorpdfstring{$v_y=v$}{vy=v}}
\[
\frac{dy}{dt}=v,\quad \frac{dx}{dt}=u(y)=\frac{2v_0}{l}y
\]
消去 $dt$：$\displaystyle dx = \frac{2v_0}{lv} y\,dy$
积分：
\[
\int_0^x dx = \int_0^y \frac{2v_0}{lv} y'\,dy' \implies \boldsymbol{x=\frac{v_0}{lv}y^2}
\]
去程轨迹为抛物线。到达 $y=l/4$：
\[
x_1 = \frac{v_0}{lv}\left(\frac{l}{4}\right)^2 = \frac{v_0 l}{16v}
\]
\subsubsection{第二阶段：返程 \texorpdfstring{$y=l/4 \to 0$}{y=l/4到0}，船垂直返回 \texorpdfstring{$v_y=-0.5v$}{vy=-0.5v}}
\[
\frac{dy}{dt}=-\frac{v}{2},\quad \frac{dx}{dt}=\frac{2v_0}{l}y
\]
\[
dt = \frac{dy}{-v/2}=-\frac{2dy}{v}
\]
代入x方向微分：
\[
dx = \frac{2v_0}{l}y \cdot \left(-\frac{2dy}{v}\right) = -\frac{4v_0}{lv} y\,dy
\]
积分上下限：$x:x_1 \to x_\text{总},\ y:l/4 \to 0$
\[
\int_{x_1}^{x_\text{总}} dx = \int_{l/4}^{0} -\frac{4v_0}{lv}\,y\,dy
\]
右侧积分：
\[
\int_{l/4}^{0} -\frac{4v_0}{lv}y\,dy
= -\frac{4v_0}{lv}\cdot\left.\frac{y^2}{2}\right|_{l/4}^{0}
= -\frac{2v_0}{lv}\left(0-\frac{l^2}{16}\right) = \frac{v_0 l}{8v}
\]
\[
x_\text{总} - x_1 = \frac{v_0 l}{8v} \implies x_\text{总} = \frac{v_0 l}{16v}+\frac{2v_0 l}{16v}
\]
\[
\boxed{x_\text{总} = \frac{3v_0 l}{16v}}
\]
即船回到原岸时，在出发点下游 $\boldsymbol{\dfrac{3v_0 l}{16v}}$ 处。
\end{solution}
\subsection{飞机空投救生舱平抛运动}
\begin{center}
\begin{tikzpicture}[scale=1]
% 竖直高度h
\draw (0,5) -- (0,0) node[midway,left]{$h=500\ \mathrm{m}$};
\draw[->] (0,5) -- (2,5) node[above]{$v_0$};
% 平抛抛物线
\draw[thick,domain=0:4] plot (\x,{5 - 0.3125*\x*\x}) node[midway,above]{轨道};
% 视线斜线
\draw[dashed] (0,5) -- (4,0) node[midway,right]{视线};
% 地面
\draw (-1,0) -- (5,0);
\end{tikzpicture}
\end{center}
一架营救飞机以 $198\ \mathrm{km\cdot h^{-1}}$ 的速率在 $500\ \mathrm{m}$ 高处向落水者正上方飞去。
\begin{enumerate}
    \item 释放救生舱时，驾驶员视线与竖直方向夹角多大？
    \item 救生舱落地时速度大小与方向。
\end{enumerate}
\begin{solution}
先单位换算：
\[
v_0 = 198\ \mathrm{km/h} = 55\ \mathrm{m/s},\quad h=500\ \mathrm{m},\ g=9.8\ \mathrm{m/s^2}
\]
\begin{enumerate}
    \item \textbf{下落时间}
竖直自由下落：
\[
h = \frac12 g t^2 \implies t=\sqrt{\frac{2h}{g}}
=\sqrt{\frac{1000}{9.8}}\approx 10.10\ \mathrm{s}
\]
水平射程：
\[
x = v_0 t = 55 \times 10.10 \approx 555.5\ \mathrm{m}
\]
设竖直偏角为 $\theta$：
\[
\tan\theta = \frac{x}{h} = \frac{555.5}{500}\approx 1.111
\]
\[
\boxed{\theta \approx 48.0^\circ}
\]
    \item \textbf{落地速度}
水平分量：$v_x=v_0=55\ \mathrm{m/s}$
竖直分量：
\[
v_y = gt = \sqrt{2gh} = \sqrt{2\times9.8\times500}\approx98.99\ \mathrm{m/s}
\]
合速度大小：
\[
v=\sqrt{v_x^2+v_y^2}
=\sqrt{55^2+98.99^2}\approx \boxed{113.3\ \mathrm{m/s}}
\]
速度与竖直方向夹角 $\alpha$：
\[
\tan\alpha = \frac{v_x}{v_y} \approx \frac{55}{98.99} \implies \alpha\approx 29.1^\circ
\]
\end{enumerate}
\end{solution}

% !TeX root = main.tex
% 第二章 质点动力学习题解答 子文件
% solution 环境在主文档main.tex中定义，本子文件不再重复定义
\section{质点动力学习题解答}
\subsection{喷气式飞机引擎推力}
某喷气式飞机以 $200\ \mathrm{m\cdot s^{-1}}$ 的速度在空中飞行，引擎每秒吸入 $50\ \mathrm{kg}$ 的空气与 $2\ \mathrm{kg\cdot s^{-1}}$ 的燃料混合燃烧，燃烧后的气体相对飞机以 $400\ \mathrm{m\cdot s^{-1}}$ 的速度向后喷出，试求该喷气式飞机引擎的推力。
\begin{solution}
取匀速飞行的飞机为惯性参考系，取 $\Delta t=1\ \mathrm{s}$。
$\Delta t$ 内吸入空气质量 $m_a=50\ \mathrm{kg}$，消耗燃料质量 $m_f=2\ \mathrm{kg}$。
空气进入引擎前相对飞机向前：$v_1=200\ \mathrm{m\cdot s^{-1}}$；
燃烧后气体相对飞机向后喷出：$v_2=-400\ \mathrm{m\cdot s^{-1}}$。
对气体应用动量定理，引擎对气体的冲量等于气体动量变化：
\[
F \Delta t = (m_a+m_f)v_2 - m_a v_1
\]
代入 $\Delta t=1\ \mathrm{s}$：
\[
F = (50+2)\cdot(-400) - 50 \cdot 200 = -30800\ \mathrm{N}
\]
$F$ 为引擎对气体的作用力，负号表示方向向后；根据牛顿第三定律，气体对引擎的推力向前：
\[
\boxed{F_\text{推}=30800\ \mathrm{N}=3.08\times 10^4\ \mathrm{N}}
\]
\end{solution}
\subsection{挖去方块平板质心（负质量法）}
一块边长 $L=6\ \mathrm{m}$ 的均匀方形平板，其一边被剪去边长 $d=2\ \mathrm{m}$ 的一块。被剪去一块的中心位于 $x=2\ \mathrm{m},\ y=0$ 处，完整平板中心在 $x=0,\ y=0$ 处，求平板剩余部分的质心位置。
\begin{solution}
设面密度为 $\sigma$。
完整大板质量：$M=\sigma L^2$；挖去小块质量：$m=\sigma d^2$。
负质量法：剩余物体 = 完整大板 + 负质量小块
\[
x_c = \frac{M\cdot x_M - m\cdot x_m}{M-m},\quad y_c = \frac{M\cdot y_M - m\cdot y_m}{M-m}
\]
已知 $x_M=0,\ y_M=0;\ x_m=2,\ y_m=0$
\[
x_c=\frac{\sigma L^2 \cdot 0 - \sigma d^2 \cdot 2}{\sigma(L^2-d^2)}
=\frac{-2d^2}{L^2-d^2}
\]
代入 $L=6,\ d=2$：
\[
x_c=\frac{-2\times 4}{36-4}=\frac{-8}{32}=-0.25\ \mathrm{m},\quad y_c=0
\]
\[
\boxed{(x_c,\,y_c)=(-0.25\ \mathrm{m},\ 0)}
\]
\end{solution}
\subsection{蓄水池抽水做功与抽水时间}
一长方体蓄水池，底面积 $S=50\ \mathrm{m^2}$，贮水深度 $h_1=1.5\ \mathrm{m}$。假定水平面距离地面的高度为 $h_2=5\ \mathrm{m}$。问要将这池水全部抽到地面上来，抽水机需做功多少？若抽水机的效率为 $80\%$，输入功率为 $P=35\ \mathrm{kW}$，则抽完这池水需要多长时间？
\begin{solution}
水总质量：$m=\rho S h_1$，水的重心初始高度：$h_c=h_2-\dfrac{h_1}{2}$。
抽水机有用功等于水重力势能增量：
\[
W_\text{有}=mg h_c=\rho S h_1 g \left(h_2-\frac{h_1}{2}\right)
\]
代入 $\rho=1000\ \mathrm{kg\cdot m^{-3}},\ g=9.8\ \mathrm{m\cdot s^{-2}},\ S=50\ \mathrm{m^2},\ h_1=1.5\ \mathrm{m},\ h_2=5\ \mathrm{m}$
\[
W_\text{有}=1000\times 50 \times 1.5 \times 9.8 \times \left(5-\frac{1.5}{2}\right)=3.18375\times 10^6\ \mathrm{J}
\]
效率定义 $\eta=\dfrac{W_\text{有}}{W_\text{总}}$，总功 $W_\text{总}=Pt$，故
\[
t=\frac{W_\text{有}}{\eta P}
\]
代入 $\eta=0.8,\ P=35\times10^3\ \mathrm{W}$
\[
t=\frac{3.18375\times 10^6}{0.8\times 35\times 10^3}\approx 113.7\ \mathrm{s}
\]
\[
\boxed{W_\text{有}=3.18\times 10^6\ \mathrm{J},\quad t\approx 114\ \mathrm{s}}
\]
\end{solution}
\subsection{变力拉力做功（积分法）}
用一变力 $F_\mathrm{T}=10\sin\alpha\ (\mathrm{N})$ 通过轻绳和轻滑轮将一物体从 A 处（$\alpha_1=30^\circ$）拉到 B 处（$\alpha_2=60^\circ$）。设 $h=2\ \mathrm{m}$，用积分法求拉力 $F_\mathrm{T}$ 在此过程中对物体所做的功。
\begin{solution}
几何关系：绳长 $l=\dfrac{h}{\sin\alpha}$，物体移动对应绳长缩短
\[
dl = -\frac{h \cos\alpha}{\sin^2\alpha} d\alpha
\]
拉力元功 $dW=F_T(-dl)$（$dl$ 为绳长减少，取绳缩短为正）：
\[
dW=F_T (-dl)=10\sin\alpha \cdot \frac{h \cos\alpha}{\sin^2\alpha} d\alpha=\frac{10h \cos\alpha}{\sin\alpha}d\alpha
\]
积分区间 $\alpha\colon 30^\circ \to 60^\circ$，即 $\pi/6 \to \pi/3$：
\[
W=\int_{\pi/6}^{\pi/3}\frac{10h \cos\alpha}{\sin\alpha}d\alpha
=10h \int_{\pi/6}^{\pi/3}\frac{d(\sin\alpha)}{\sin\alpha}
=10h \ln\sin\alpha \bigg|_{\pi/6}^{\pi/3}
\]
代入 $h=2\ \mathrm{m}$：
\[
W=20\left(\ln\frac{\sqrt{3}}{2}-\ln\frac12\right)=20\ln\sqrt{3}=10\ln 3 \approx 10.99\ \mathrm{J}
\]
\[
\boxed{W=10\ln 3 \approx 11.0\ \mathrm{J}}
\]
\end{solution}
\subsection{蒸汽打桩机（变阻力+完全非弹性碰撞）}
铁塔高 $40\ \mathrm{m}$，锤的质量 $M=10\ \mathrm{t}$。现将长达 $38.5\ \mathrm{m}$ 的钢筋混凝土桩打入地层，已知桩的质量为 $m=24\ \mathrm{t}$，横截面积为 $0.25\ \mathrm{m^2}$ 的正方形，桩的侧面单位面积所受泥土的阻力为 $F_r=2.65\times 10^4\ \mathrm{N\cdot m^{-2}}$。
\begin{enumerate}
    \item 靠桩自重能下沉多深？
    \item 桩稳定后，把锤提高 $1\ \mathrm{m}$，然后让锤自由下落而击桩，假定锤与桩发生完全非弹性碰撞，一锤能打下多深？
\end{enumerate}
\begin{solution}
桩截面正方形边长 $a=\sqrt{0.25}=0.5\ \mathrm{m}$，周长 $C=4a=2\ \mathrm{m}$；下沉深度 $s$ 时侧面积 $A=C s$，阻力 $f(s)=F_r \cdot C s$。
\paragraph{(1) 自重下沉深度 $s_0$}
动能定理，初末动能均为0：重力做功等于克服变阻力做功
\[
mg s_0 - \int_{0}^{s_0} F_r C s \,ds = 0
\]
\[
mg s_0 = \frac12 F_r C s_0^2 \implies s_0=\frac{2mg}{F_r C}
\]
代入 $m=24\times10^3\ \mathrm{kg},\ g=9.8\ \mathrm{m\cdot s^{-2}},\ F_r=2.65\times10^4\ \mathrm{N\cdot m^{-2}},\ C=2\ \mathrm{m}$
\[
s_0=\frac{2\times 24\times10^3\times9.8}{2.65\times10^4 \times 2}\approx 8.88\ \mathrm{m}
\]
\paragraph{(2) 锤下落撞击，一次打入深度 $\Delta s$}
锤自由下落 $h=1\ \mathrm{m}$：$v_0=\sqrt{2gh}$。
完全非弹性碰撞动量守恒：$M v_0=(M+m) v_\text{共}$
\[
v_\text{共}=\frac{M}{M+m}\sqrt{2gh}
\]
碰撞后下沉 $\Delta s$，动能定理：
\[
(M+m)g\Delta s - \int_{s_0}^{s_0+\Delta s} F_r C s\,ds = 0-\frac12(M+m)v_\text{共}^2
\]
\[
(M+m)g\Delta s - \frac12 F_r C \big[(s_0+\Delta s)^2-s_0^2\big] = -\frac{M^2 g h}{M+m}
\]
展开 $(s_0+\Delta s)^2-s_0^2=2s_0\Delta s+(\Delta s)^2$，代入数值得到关于 $\Delta s$ 的一元二次方程，取物理正根：
\[
\boxed{s_0\approx 8.88\ \mathrm{m},\quad \Delta s\approx 0.20\ \mathrm{m}}
\]
\end{solution}
\subsection{旋转匀质绳的内部张力}
一条质量为 $m$ 且分布均匀的绳子，长度为 $L$，一端拴在转轴上，并以恒定角速度 $\omega$ 在水平面上旋转，设转动过程中绳子始终伸直，且忽略重力与空气阻力，求距转轴为 $r$（$r<L$）处绳中的张力。
\begin{solution}
线密度 $\lambda = \dfrac{m}{L}$。
取距离转轴 $r$ 到 $L$ 的一段绳元作为研究对象。在位置 $x$ 处取微元 $dx$，微元质量 $dm=\lambda dx$，微元做圆周运动向心力 $dT = dm \cdot \omega^2 x$。
$r$ 处张力提供外侧全部绳段的向心力：
\[
T(r) = \int_{r}^{L} \lambda \omega^2 x \,dx
\]
\[
T(r) = \lambda \omega^2 \cdot \frac{x^2}{2}\bigg|_{r}^{L}
= \frac{m}{2L}\omega^2\left(L^2-r^2\right)
\]
\[
\boxed{T(r) = \frac{m\omega^2}{2L}\left(L^2-r^2\right)}
\]
\end{solution}
\subsection{匀质细棒浸入液体下落（动能定理）}
在密度为 $\rho_1$ 的液体上方悬一长为 $l$，密度为 $\rho_2$ 的均匀细棒，棒的 $B$ 端刚好和液面接触。今剪断细绳，设棒只在浮力和重力作用下下沉，
\begin{enumerate}
  \item 棒刚好全部浸入液体时的速度；
  \item 若 $\rho_2<\rho_1/2$ 时，求棒浸入液体的最大深度；
  \item 棒下落过程中所能达到的最大速率。
\end{enumerate}
\begin{solution}
设棒横截面积为 $S$，总质量 $m=\rho_2 S l$。
\begin{enumerate}
\item \textbf{刚好全部浸入时速度（下落距离 $l$）}
重力做功：$W_G = mg \cdot \dfrac{l}{2}$（重心下移 $l/2$）。
浸入深度 $x$ 时浮力 $F_\text{浮}(x)=\rho_1 g S x$，浮力做负功：
\[
W_\text{浮} = -\int_{0}^{l} \rho_1 g S x \,dx = -\frac12 \rho_1 g S l^2
\]
动能定理：$W_G + W_\text{浮} = \dfrac12 m v_1^2 - 0$
\[
\rho_2 S l g \cdot \frac{l}{2} - \frac12 \rho_1 g S l^2
= \frac12 \rho_2 S l \,v_1^2
\]
约去公共项整理：
\[
\rho_2 g l^2 - \rho_1 g l^2 = \rho_2 l\, v_1^2
\Rightarrow
v_1=\sqrt{\frac{(\rho_2-\rho_1)gl^2}{\rho_2 l}}
=\sqrt{\frac{(\rho_2-\rho_1)gl}{\rho_2}}
\]
注意：若 $\rho_2<\rho_1$，根号内为负，物理含义：棒不可能完全浸入；本题模型仅数学形式如上。
\[
\boxed{v_1 = \sqrt{\frac{(\rho_2-\rho_1)gl}{\rho_2}}}
\]
\item \textbf{最大浸入深度 $h$（末速度为0，$\rho_2<\rho_1/2$）}
全程重力做功：$W_G = mg \cdot \dfrac{h}{2}$。
浮力分两段：$0\to l$ 变浮力；$l\to h$ 完全浸入，恒浮力 $F_\text{浮}=\rho_1 g S l$。
\[
W_\text{浮} = -\int_{0}^{l}\rho_1 g S x dx - \int_{l}^{h}\rho_1 g S l \,dx
= -\frac12 \rho_1 g S l^2 - \rho_1 g S l (h-l)
\]
动能定理：$W_G + W_\text{浮}=0$
\[
\frac12 \rho_2 S l g h - \rho_1 g S l \left(h-\frac{l}{2}\right)=0
\]
约去公因子：
\[
\frac12 \rho_2 h - \rho_1\left(h-\frac{l}{2}\right)=0
\implies h = \frac{\rho_1 l}{2(\rho_1-\rho_2)}
\]
\[
\boxed{h_\text{max} = \frac{\rho_1 l}{2(\rho_1-\rho_2)}}
\]
\item \textbf{最大速率（合力为0，加速度 $a=0$）}
设浸入深度 $x_0$ 受力平衡：$mg = \rho_1 g S x_0$
\[
\rho_2 S l g = \rho_1 g S x_0 \implies x_0 = \frac{\rho_2}{\rho_1}l
\]
下落 $x_0$ 过程动能定理：
\[
mg\cdot \frac{x_0}{2} - \int_{0}^{x_0}\rho_1 g S x dx = \frac12 m v_\text{max}^2
\]
代入 $m=\rho_2 S l,\ x_0=\dfrac{\rho_2}{\rho_1}l$，化简：
\[
\boxed{v_\text{max} = \sqrt{\frac{\rho_2 g l^2}{\rho_1}\left(2-\frac{\rho_2}{\rho_1}\right)}}
\]
\end{enumerate}
\end{solution}
\subsection{F‑t图像 动量定理求7s末速度、坐标}
一质点沿 $x$ 轴运动，其所受的力如题图所示，设 $t=0$ 时，$v_0=5\ \mathrm{m\cdot s^{-1}}$，$x_0=2\ \mathrm{m}$，质点质量 $m=1\ \mathrm{kg}$，试求该质点 $7\ \mathrm{s}$ 末的速度和坐标。
\begin{solution}
动量定理：$\displaystyle \int_{0}^{t} F\,dt = m(v-v_0)$，$F$‑$t$ 图像面积为冲量。
分段力函数：
\[
\begin{cases}
0\le t\le 5\colon & F_1 = 2t \\
5< t\le 7\colon & F_2 = 35 -5t
\end{cases}
\]
\textbf{0~5 s：}
\[
I_1=\int_0^5 2t\,dt = \big[t^2\big]_0^5=25\ \mathrm{N\cdot s}
\]
\[
a_1=\frac{F_1}{m}=2t,\quad
v_1(t) = v_0+\int_0^t 2t'dt' = 5 + t^2
\]
\[
x_1(t)=x_0+\int_0^t v_1(t')dt' = 2 +5t+\frac{t^3}{3}
\]
$t=5\ \mathrm{s}$：
\[
v_5=5+5^2=30\ \mathrm{m\cdot s^{-1}},\quad
x_5=2+25+\frac{125}{3}=\frac{206}{3}\ \mathrm{m}
\]
\textbf{5~7 s：}
\[
I_2=\int_{5}^{7}(35-5t')dt'
=\left[35t'-\frac52 t'^2\right]_{5}^{7}
=(245-122.5)-(175-62.5)=10\ \mathrm{N\cdot s}
\]
总冲量 $I=I_1+I_2=25+10=35\ \mathrm{N\cdot s}$。
\[
v_7 = v_0+\frac{I}{m}=5+35=40\ \mathrm{m\cdot s^{-1}}
\]
速度函数 $t\in[5,7]$：
\[
v_2(t)=v_5+\int_{5}^{t}(35-5t')dt'
=30+\left(35(t-5)-\frac52(t^2-25)\right)
\]
位置：
\[
x_7 = x_5+\int_{5}^{7} v_2(t')dt' =\frac{418}{3}\approx 139.33\ \mathrm{m}
\]
\[
\boxed{v_7=40\ \mathrm{m\cdot s^{-1}},\quad x_7=\frac{418}{3}\approx 139.33\ \mathrm{m}}
\]
\end{solution}
\subsection{连接体，绳子绷紧后C开始运动（动量守恒）}
三物体质量都为 $m$，B 和 C 紧靠在一起放在光滑的水平面上，二者间连有一段长为 $0.4\ \mathrm{m}$ 的绳子，B 的另一侧的绳子通过滑轮和物体 A 相连，若滑轮和绳子质量不计，滑轮轴上摩擦也不计，绳长一定，问 A，B 开始运动后，经多少时间 C 也开始运动？运动速度为多少？
\begin{solution}
第一阶段：A、B 一起匀加速，C 静止，位移 $s=0.4\ \mathrm{m}$。
对 A、B 整体：$mg = 2ma \implies a=\dfrac{g}{2}$。
匀加速 $s=\dfrac12 a t^2$：
\[
0.4 = \frac12 \cdot \frac{g}{2} \cdot t^2
\implies t=\sqrt{\frac{1.6}{g}}
\]
取 $g=9.8\ \mathrm{m\cdot s^{-2}}$：
\[
\boxed{t \approx 0.404\ \mathrm{s}}
\]
绳子绷紧前瞬间 A、B 速度：
\[
v_0 = a t = \sqrt{2 a s} = \sqrt{2\cdot \frac{g}{2}\cdot 0.4}=\sqrt{0.4g}
\]
绷紧瞬间内力远大于外力，ABC 系统水平动量守恒：
\[
2m v_0 = 3m v
\implies v=\frac23 v_0
\]
\[
\boxed{v = \frac{2}{3}\sqrt{0.4g}\approx 1.32\ \mathrm{m\cdot s^{-1}}}
\]
\end{solution}
\subsection{浮吊水平移动（水平方向质心守恒）}
浮吊，质量为 $m'=2\times 10^4\ \mathrm{kg}$，由岸上吊起质量 $m=2\times 10^3\ \mathrm{kg}$ 的重物后，再将吊杆与竖直方向的夹角 $\theta$ 由 $60^\circ$ 转到 $30^\circ$。设杆长 $l=8\ \mathrm{m}$，水的阻力和杆重忽略不计。求浮吊在水平方向移动的距离。
\begin{solution}
系统水平方向不受外力，水平质心坐标保持不变。
设浮吊相对岸向岸边移动距离为 $s$（浮吊对地水平位移 $-s$）；
重物相对浮吊的水平位移大小：$l(\sin60^\circ-\sin30^\circ)$，
重物相对地面水平位移：$l(\sin60^\circ-\sin30^\circ)-s$。
质心守恒：
\[
m'(-s)+m\big[l(\sin60^\circ-\sin30^\circ)-s\big]=0
\]
\[
s(m'+m)=ml\left(\frac{\sqrt{3}}{2}-\frac12\right)
\]
\[
s = \frac{ml(\sqrt{3}-1)}{2(m'+m)}
\]
代入数值：
\[
s=\frac{2\times10^3 \times 8 (\sqrt{3}-1)}{2(2\times10^4+2\times10^3)}
=\frac{8(\sqrt{3}-1)}{22}\approx 0.266\ \mathrm{m}
\]
\[
\boxed{s\approx 0.266\ \mathrm{m}，浮吊向岸边移动约0.266\ \mathrm{m}}
\]
\end{solution}

\input{5.刚体力学基础.tex}
\input{6.流体力学.tex}
\input{7.狭义相对论.tex}
% ========== 第二篇 电磁学 ==========
\part{第二篇 电磁学}
\input{8.电荷与电场.tex}
\input{9.电流与磁场.tex}
\input{10.电磁场与麦克斯韦电磁场方程组.tex}
% ========== 第三篇 振动与波动 ==========
\part{第三篇 振动与波动}
\input{11.振动学基础.tex}
\input{12.波动学基础.tex}
\end{document}
